\documentclass{article}

\usepackage{arxiv}

\usepackage[spanish]{babel}
\usepackage[utf8]{inputenc}
\usepackage[T1]{fontenc}
\usepackage{amsmath}
\usepackage{amssymb}
\usepackage{hyperref}
\usepackage{cleveref}
\usepackage{url}
\usepackage{booktabs}
\usepackage{amsfonts}
\usepackage{nicefrac}
\usepackage{microtype}
\usepackage{graphicx}
\usepackage{float}
\usepackage{subcaption}
\crefname{table}{cuadro}{cuadros}
\Crefname{table}{Cuadro}{Cuadros}
\graphicspath{{./images/}{./images/graficos/}}


\title{TP 3 - SIMULACIÓN DE UN MODELO MM1 E INVENTARIO}


\author{
    Alvaro Tomasino \\
    Ingeniería en Sistemas de Información\\
    Universidad Tecnológica Nacional - FRRO\\
    Zeballos 1341, S2000, Argentina\\
    \texttt{atomasino08@gmail.com} \\
\And
    Juan Bautista Martinelli \\
    Ingeniería en Sistemas de Información\\
    Universidad Tecnológica Nacional - FRRO\\
    Zeballos 1341, S2000, Argentina\\
    \texttt{martinellij4@gmail.com} \\
\And
    Tomás Aguilera \\
    Ingeniería en Sistemas de Información\\
    Universidad Tecnológica Nacional - FRRO\\
    Zeballos 1341, S2000, Argentina\\
    \texttt{tomas.aguilera.27@gmail.com} \\
\And
    Melina Aronson \\
    Ingeniería en Sistemas de Información\\
    Universidad Tecnológica Nacional - FRRO\\
    Zeballos 1341, S2000, Argentina\\
    \texttt{melinaaronson@gmail.com} \\
}




\begin{document}

\maketitle




\begin{abstract}
El objetivo de este informe es medir qué tan bien funcionan el sistema de colas M/M/1 y el modelo de inventario bajo distintas situaciones, comparando los cálculos matemáticos con programas en Python y modelos en AnyLogic.
\end{abstract}

\keywords{Simulación \and Modelo MM1 \and Modelo de Inventario \and AnyLogic}




\section{Introducción}

En cualquier negocio, manejar bien las filas de espera y el stock de productos es fundamental para no perder dinero ni tiempo. La teoría de colas nos da herramientas para entender sistemas donde la gente llega y espera ser atendida. El modelo M/M/1 permite predecir, por ejemplo, cuánto tiempo va a estar un cliente en una fila antes de irse.

Por otro lado, los modelos de inventario sirven para controlar el stock, buscando un equilibrio entre el costo de pedir mercadería, mantenerla guardada y el problema de quedarse sin nada para vender.

Aunque existen fórmulas para resolver estos problemas, la simulación computacional es mejor para ver cómo se comportan estos sistemas en el tiempo y con cambios inesperados. En este informe, vamos a comparar los resultados teóricos con los obtenidos en Python y AnyLogic. La idea es probar distintos escenarios para ver si las herramientas de simulación coinciden con la matemática.




\section{Marco teórico y fórmulas}
\label{sec:marco}

\subsection{Modelo M/M/1}

\subsection{Modelo de Inventario}





\section{Metodología}
\label{sec:metodologia}

\subsection{Python}

\subsubsection{Modelo de Colas MM1}

\subsubsection{Modelo de Inventario}

\subsection{AnyLogic}

\subsubsection{Implementación del Modelo M/M/1 y M/M/1/K}

\subsubsection{Implementación del Modelo de Inventario}

\subsubsection{IU y experimentación}



\subsection{Parámetros elegidos para el modelo de Inventario}






\section{Resultados}
\label{sec:resultados}

\subsection{Análisis del Modelo M/M/1}

\subsubsection{Evolución Temporal y Estabilización}

\subsubsection{Sensibilidad ante la tasa de arribo}

\subsubsection{Análisis de cola finita (M/M/1/K) y Denegación de servicio}

\subsubsection{Probabilidad de encontrar n clientes en cola}


\subsection{Análisis del Modelo de Inventario}

\subsubsection{Comparativa entre valor teórico, Python y AnyLogic}

\section{Discusión}
\label{sec:discusion}


\section{Conclusiones}
\label{sec:conclusiones}



\nocite{*}

\bibliographystyle{unsrturl}
\bibliography{references}

\end{document}